\documentclass[10pt,letterpaper]{article}
\usepackage[letterpaper,margin=1.6cm]{geometry}
\usepackage[T1]{fontenc}
\usepackage[utf8]{inputenc}
\usepackage{newtxtext}
\usepackage{xcolor}
\usepackage[hidelinks]{hyperref}
\usepackage{enumitem}
\pagestyle{empty}
\setlength{\parindent}{0pt}
\setlength{\parskip}{2pt}
\linespread{0.96}
\setlist[itemize]{leftmargin=5mm,itemsep=0pt,topsep=1pt,parsep=0pt}
\definecolor{ieeeblue}{RGB}{0,84,166}
\newcommand{\clhead}[1]{\textcolor{ieeeblue}{\textbf{#1}}}

\begin{document}

\begin{flushright}
\textbf{St\'ephane Ga\"el R. Ekodeck}\\
Department of Computer Science, LIRIMA, Team GRIMCAPE\\
University of Yaound\'e I, P.O.\ Box 812, Yaound\'e, Cameroon\\
\href{mailto:stephane-gael.ekodeck@facsciences-uy1.cm}%
{stephane-gael.ekodeck@facsciences-uy1.cm}\\
June 2026
\end{flushright}

The Editor-in-Chief\\
\textit{IEEE Transactions on Multimedia}\\
IEEE Signal Processing Society\\

Dear Editor-in-Chief,

We respectfully submit our manuscript, \textbf{``NARCIS: Authenticated
Neural Coverless Image Signaling with Attack-Qualified Codebooks and
Error Correction,''} for consideration as an original research article in
\textit{IEEE Transactions on Multimedia} (TMM).

\clhead{Scope alignment.}
NARCIS addresses a core multimedia communications problem: how to
transmit authenticated messages by selecting and sending unchanged
natural images from a shared index, rather than by modifying a carrier.
The technical pipeline is a multimedia processing contribution ---
a compact self-supervised encoder that maps images to channel-invariant
descriptors, a quantile codebook that converts the descriptor space into
balanced message symbols, and a multi-attack qualification stage that
retains only images whose labels remain stable under channel
distortions. These components sit squarely within TMM's scope on
multimedia content analysis, multimedia security, and image-based
communications. The detectability evaluation (three detector classes,
five partitions per dataset, calibration and holdout conditions) also
directly addresses the multimedia security analysis that TMM regularly
publishes.

\clhead{Key technical contributions.}
\begin{itemize}
\item \textbf{Feasible-codebook operating point (Contribution~1).}
  Existing coverless systems report nominal capacity or average symbol
  accuracy, neither of which verifies that a given message can actually
  be transmitted from a finite, qualified index. NARCIS introduces the
  \emph{feasible operating point}: the largest symbol alphabet for which
  every attack-qualified bucket satisfies the multiplicity required by
  every protected coded message. Over 1{,}575 complete-message trials on
  BOSSBase~1.01 (grayscale) and Caltech-101 (RGB), NARCIS recovered
  100\% of authenticated plaintexts under 21 channel conditions,
  including eight unseen holdout strengths.
\item \textbf{End-to-end coverless pipeline (Contribution~2).}
  A self-supervised encoder (variance--invariance--covariance objective)
  produces 64-dimensional normalised descriptors. A quantile partition
  of the first principal component yields balanced symbol classes.
  Multi-attack qualification (12 calibration transforms) removes
  geometrically unstable images; a calibration-locked projection search
  maintains 210/210 recoveries while raising the minimum stable bucket
  from 92 to~161 (BOSSBase seed~11). AES-GCM authenticates payloads
  and session metadata; a keyed Gray permutation limits symbol exposure;
  Reed--Solomon coding corrects residual label errors. Net throughput
  scales from 0.184 to 1.213~bits per cover for 8- to 128-byte payloads.
\item \textbf{Detectability and ablation evaluation (Contribution~3).}
  Selection bias is quantified with three detector classes (global
  logistic regression on dense residuals, residual-feature ExtraTrees,
  selection CNN) over five partitions each of BOSSBase and Caltech-101.
  BOSSBase AUCs include chance. Caltech-101 classical detectors identify
  a weak but reproducible selection shift (global logistic AUC~0.531,
  95\%~CI~[0.504, 0.557]; ExtraTrees~0.533, [0.504, 0.562]); 20
  repeated CNN fits on the most sensitive partition confirm a
  partition-specific signal (mean~AUC~0.533, 95\%~CI~[0.517, 0.548]),
  reported transparently rather than dismissed. Removing attack
  qualification on BOSSBase seed~11 reduces complete-message success
  from 210/210 to 142/210~(67.6\%), establishing qualification as the
  dominant robustness mechanism.
\end{itemize}

\clhead{Novelty and positioning.}
Selection-based coverless systems have been evaluated at the descriptor
level --- average symbol accuracy, nominal capacity --- but have not
verified message-level feasibility over a finite qualified index, nor
integrated payload authentication, replay control, and error correction
into a single testable system. TMM has published related work in
learned descriptors for coverless steganography (Ren--Wu~2024;
Guo--Ping~2026) and in multimedia security evaluation. NARCIS
complements this body of work by introducing the feasible-codebook
criterion as the primary evaluation unit and providing systematic
cross-dataset, multi-detector detectability evidence.

On capacity, the paper distinguishes two metrics that are not
interchangeable: \emph{nominal symbol capacity} ($\log_2 K$ bpc,
measured at the symbol level before any protocol overhead) and
\emph{net authenticated throughput} (plaintext bpc after AES-GCM,
Reed-Solomon parity, and framing). NARCIS achieves 3--4~bpc nominal at
$K{=}8$--$16$ and 0.184--1.213~bpc net. Prior systems reporting
10--15~bpc measure nominal capacity without authentication overhead;
adding a comparable protocol stack would reduce their net throughput by
the same constant cost. The paper reports both metrics explicitly and
does not conflate them.

\clhead{Author declarations.}
\begin{itemize}
\item The manuscript is original, has not been published, and is not
  under consideration elsewhere.
\item All four authors approved the final manuscript and author order.
  WYSAWIS (cited in the related work) shares authors with this
  submission and is explicitly identified as such in the paper; it is
  used as published qualitative context, not as an experimental
  baseline.
\item No competing interests and no external funding.
\item No human subjects, personal data, or animals are involved.
  Implementation, fixed seeds, experimental results, and manuscript
  sources are available at
  \href{https://github.com/EkodeckStephane/NARCIS}%
  {github.com/EkodeckStephane/NARCIS}.
\end{itemize}

We believe NARCIS will interest TMM readers working on multimedia
steganography, learned image representations, and multimedia security
evaluation. We welcome the opportunity to provide any additional
information the Associate Editor may require.\\[4pt]

Yours sincerely,\\[6pt]

\textbf{St\'ephane Ga\"el R. Ekodeck} (Corresponding Author)\\
ORCID: 0000-0002-8094-8832\\
\textbf{Co-authors:} Vivien Beyala Kamgang (ORCID 0000-0002-3644-1866);
Serge Alain Eb\'el\'e (ORCID 0000-0001-7036-7068);
Julius Marcellin Nkenlifack, Member IEEE (ORCID 0000-0003-3322-4687).

\end{document}
